% =========================================================
% CONCLUSIONES — versión RUCP (~600 palabras)
% =========================================================

\section{Conclusiones}
\label{sec:conclusiones}

Este estudio aplicó inferencia ecológica bayesiana sobre 7{,}271 circuitos electorales para reconstruir las transferencias de votos en el balotaje presidencial uruguayo de 2024 y su comparación con 2019. Los resultados revelan tres hallazgos sustantivos y una implicación metodológica.

\textbf{La rebelión rural} es el hallazgo más contraintuitivo: los circuitos rurales de Uruguay exhibieron tasas de defección de Cabildo Abierto al Frente Amplio 21{,}8 pp superiores a los urbanos (60{,}8\% vs.\ 39{,}0\%), con $P(T_{\text{rural}} > T_{\text{urbano}} \mid \text{datos}) > 0{,}99$. La narrativa del conservadurismo rural, relevante en contextos de elecciones de primera vuelta, se invierte en el contexto del balotaje de 2024, probablemente como reflejo de la evaluación retrospectiva negativa del desempeño gubernamental en zonas donde el impacto de la crisis agrícola fue más directo.

\textbf{La paradoja de la coalición cohesiva derrotada}: tres de los cuatro socios de la coalición mejoraron sustancialmente su retención de votos entre 2019 y 2024 (CA: $+$42.4 pp de mejora en la dirección esperada; PC: $-$26.6 pp de defección; PN: $-$32.4 pp de defección; PI: $-$92.1 pp). Sin embargo, la coalición perdió el balotaje por $\approx$95{,}000 votos. La causa es el colapso del caudal de CA ($-$78\% entre elecciones): la coalición ingresó a la segunda vuelta de 2024 con 154{,}000 votos menos que en 2019 para transferir hacia el candidato de la coalición. \textit{La lección política es directa: la cohesión en la segunda vuelta no compensa la pérdida en la primera.} Las coaliciones deben gestionar el crecimiento o mantenimiento de su caudal en la primera vuelta, no solo la transferencia en la segunda.

\textbf{La reconfiguración geográfica}: la nula correlación entre patrones de defección departamentales de 2019 y 2024 ($r < 0.13$) indica que no existe una geografía estable de la deslealtad coalicional. Los departamentos con alta defección en 2019 no fueron los mismos en 2024, y viceversa. Esta inestabilidad hace que la planificación electoral basada en mapas históricos de fidelidad sea de valor limitado; las condiciones coyunturales locales dominan sobre las identidades geográficas persistentes.

\textbf{Implicación metodológica}: el marco bayesiano de inferencia ecológica permite aprovechar directamente la distribución posterior completa para comparaciones geográficas formales ($P(T_A > T_B \mid \text{datos})$), superando al enfoque frecuente de comparar si los intervalos de confianza se superponen. Esta ventaja del enfoque bayesiano ---documentada aquí en el contexto de la heterogeneidad urbano-rural--- es transferible a cualquier análisis de transferencias de votos en sistemas multipartidistas con elecciones de dos rondas.

Los resultados contribuyen a la literatura sobre comportamiento electoral en presidencialismos latinoamericanos, sobre divisiones urbano-rurales en democracias contemporáneas, y sobre metodología de inferencia ecológica aplicada a sistemas con múltiples partidos. Las extensiones prioritarias ---modelo jerárquico con covariables censales y triangulación con métodos alternativos--- están delineadas como trabajo futuro concreto.
